\documentclass[12pt,a4paper]{article}
\usepackage[utf8]{inputenc}
\usepackage[T1]{fontenc}
\usepackage{lmodern}
\usepackage[hungarian]{babel}
\usepackage{geometry}
\geometry{margin=2.5cm}
\usepackage{graphicx}
\usepackage{booktabs}
\usepackage{hyperref}
\usepackage{bookmark}
\usepackage{xcolor}
\usepackage{listings}
\usepackage{amsmath}
\usepackage{float}
\usepackage{enumitem}
\usepackage[hidelinks]{hyperref}

\hypersetup{
    colorlinks=true,
    linkcolor=blue!60!black,
    urlcolor=blue!60!black,
    citecolor=blue!60!black,
}

\lstset{
    basicstyle=\ttfamily\small,
    keywordstyle=\color{blue!70!black}\bfseries,
    commentstyle=\color{green!50!black}\itshape,
    stringstyle=\color{red!60!black},
    backgroundcolor=\color{gray!8},
    frame=single,
    framerule=0.4pt,
    rulecolor=\color{gray!50},
    breaklines=true,
    showstringspaces=false,
    numbers=left,
    numberstyle=\tiny\color{gray},
    language=Python,
    tabsize=4,
}

\title{
    \textbf{Ipari Képfeldolgozás Laboratóriumi Gyakorlat} \\[0.3cm]
    \Large 3.~Mérföldkő — Végleges Algoritmus \& Eredmények \\[0.3cm]
    \large Projekt: Objektumok regisztrációja 3D pontfelhőn
}
\author{}
\date{}

\begin{document}
\maketitle
\thispagestyle{empty}

\vspace{-1.5cm}

\begin{table}[H]
\centering
\begin{tabular}{ll}
\toprule
\textbf{Beadási határidő} & 2026.~május 15.~23:59 \\
\textbf{Csapat tagjai}     & Ilonczai András \\
\textbf{Dátum}             & 2026.~május \\
\bottomrule
\end{tabular}
\end{table}

% ============================================================
\section{Projekt összefoglalás és végső státusz}
% ============================================================

Az 1.~és 2.~mérföldkövek alapján a 3D pontfelhő regisztrációs pipeline teljes implementációja megtörtént és optimalizálásra került. A legújabb fejlesztések a robusztusság javítására és az interaktív vizualizáció kiegészítésére összpontosítottak, melyek az alábbi kérdéseket oldják meg:

\begin{itemize}[nosep]
    \item \textbf{Regisztrációs robusztusság:} Dinamikus paraméter hangolás a különböző pontfelhő karakterisztikák kezelésére
    \item \textbf{Vizualizáció:} PyVista-alapú interaktív 3D megjelenítő fájlválasztó funkcionalitással
    \item \textbf{Felhasználói élmény:} Intuitív terminál-alapú felület az eredmények megtekintéséhez
\end{itemize}

A végleges rendszer megbízható regisztrációt biztosít különféle 3D pontfelhő adathalmazokra, amit egy saját modellon végzett tesztelés is igazolt.

% ============================================================
\section{Végleges algoritmus}
% ============================================================

Az előző mérföldkövekben bemutatott hat-lépéses pipeline az alábbiak szerint lett finalizálva és optimalizálva:

\subsection{Pipeline architektúra — végső verzió}

\begin{table}[H]
\centering
\begin{tabular}{clp{8cm}}
\toprule
\textbf{Lépés} & \textbf{Megnevezés} & \textbf{Finalizálás \& Javítások} \\
\midrule
1 & Adatbevitel \& előfeldolgozás & Robusztus hiba-kezelés, geometriai validáció \\
2 & Jellemzők kinyerése & Stabil normális becslés, hibrid keresési paraméterek \\
3 & Durva illesztés (RANSAC) & Dinamikus max\_iter, konfidencia tuning \\
4 & Finom illesztés (ICP) & Point-to-Point és Point-to-Plane variánsok \\
5 & Minőségértékelés & Fitness, RMSE, korespondenciák száma \\
6 & Eredmény exportálás & PLY, JSON transzformációk, vizualizáció \\
\bottomrule
\end{tabular}
\caption{A pipeline végső hat feldolgozási lépése optimalizálásokkal.}
\end{table}

\subsection{Regisztrációs robusztusság javítása}

A regisztrációs algoritmus robusztusságának javítása érdekében a következő módosítások történtek:

\begin{enumerate}[nosep]
    \item \textbf{Dinamikus paraméter hangolás:} A RANSAC és ICP paraméterek az aktuális pontfelhő jellemzői alapján dinamikusan állítódnak be.
    \item \textbf{Kiterjesztett hiba-kezelés:} Geometriai validáció az előfeldolgozás során:
        \begin{itemize}
            \item Üres pontfelhőket kezel
            \item Hibás pontokat szűr
            \item Extrém kiterjedést detektál
        \end{itemize}
    \item \textbf{Optimalizált Statistical Outlier Removal:} 20 szomszéd, $2\sigma$ küszöbérték az izolált pontok eltávolítására.
    \item \textbf{Konzisztens normális orientáció:} A normálisok \texttt{orient\_normals\_consistent\_tangent\_plane(k=30)} módszerrel konzisztensek.
    \item \textbf{Adaptív ICP küszöbök:} Az ICP correspondence distance threshold dinamikusan számítódik a voxel size-tól.
\end{enumerate}

\subsection{Interaktív vizualizáció PyVista alapokkal}

Az interaktív vizualizáció teljes átalakítása PyVista alapokkal:

\begin{enumerate}[nosep]
    \item \textbf{PyVista 3D megjelenítő:} Off-screen és interaktív renderelési módok támogatása.
    \item \textbf{Interaktív vezérlés:}
        \begin{itemize}
            \item \textbf{Bal egér:} Forgatás
            \item \textbf{Jobb egér / scroll:} Zoom
            \item \textbf{Középső gomb:} Mozgatás
            \item \textbf{R billentyű:} Kamera reset
            \item \textbf{Q billentyű:} Kilépés
        \end{itemize}
    \item \textbf{Fájlválasztó interface:}
        \begin{itemize}
            \item Terminál-alapú checkbox rendszer
            \item Támogatott formátumok: számok (1,3,5), tartományok (1-4), összes (all), output merged (m)
            \item Dinamikus legenda a megjelenített fájlok szerint
        \end{itemize}
    \item \textbf{Szín-kódolás:} Nyolc előre definiált szín az egyidejűleg megjeleníthető pontfelhőkhöz.
    \item \textbf{Valós idejű hibajelzés:} Korrupt vagy nem betölthető fájlok kezelése.
\end{enumerate}

% ============================================================
\section{Technikai megvalósítás}
% ============================================================

\subsection{Adatbevitel és robusztus előfeldolgozás}

Az előfeldolgozás során a következő robusztusságjavítások történtek:

\begin{itemize}[nosep]
    \item \textbf{Betöltés és validáció:} Open3D pontfelhő betöltése és alapvető ellenőrzés
    \item \textbf{Hibás pontok eltávolítása:} Non-finite pontok szűrése, Statistical Outlier Removal (20 szomszéd, $2\sigma$ küszöbérték)
    \item \textbf{Geometriai validáció:} Az objektum kiterjedésének ellenőrzése (üres vagy extrém nagy pontfelhő detektálása)
    \item \textbf{Voxel downsampling:} A pontfelhő ritkítása a megadott voxel méret szerint
\end{itemize}

\subsection{Jellemző kinyerés — normálisok és FPFH}

Az optimalizált jellemző kinyerés a következő lépésekből áll:

\begin{itemize}[nosep]
    \item \textbf{Normális becslés:} PCA-alapú módszer hibrid keresési paraméterekkel (sugár: 2 × VOXEL\_SIZE, max 30 szomszéd)
    \item \textbf{Konzisztens orientáció:} Normálisok orientálása konzisztens tangent plane módszerrel
    \item \textbf{FPFH deszkriptorok:} Fast Point Feature Histograms számítása (33 dimenziós), keresési sugár: 5 × VOXEL\_SIZE
\end{itemize}

\subsection{Durva illesztés — Dinamikus RANSAC}

A RANSAC algoritmus dinamikus paraméterezésű verzióval működik:

\begin{table}[H]
\centering
\begin{tabular}{ll}
\toprule
\textbf{Paraméter} & \textbf{Végső érték / Szabály} \\
\midrule
Távolság küszöb & $1{,}5 \times \text{VOXEL\_SIZE}$ \\
Maximális iteráció & 200\,000 -- 200\,000 (dinamikus) \\
Konfidencia & 0{,}9999 (99{,}99\%) \\
Minta pontok & 3 \\
Edge length checker & $\geq 0{,}9$ \\
\bottomrule
\end{tabular}
\caption{RANSAC végső paraméterezése.}
\end{table}

\subsection{Finom illesztés — Adaptív ICP}

Az ICP két variánsban áll rendelkezésre, dinamikus küszöbökkel:

\paragraph{Point-to-Point ICP:}
\[
E(\mathbf{R}, \mathbf{t}) = \sum_{i=1}^{N} \| \mathbf{R} \mathbf{p}_i + \mathbf{t} - \mathbf{q}_i \|^2
\]

\paragraph{Point-to-Plane ICP:}
\[
E(\mathbf{R}, \mathbf{t}) = \sum_{i=1}^{N} \left[ (\mathbf{R} \mathbf{p}_i + \mathbf{t} - \mathbf{q}_i) \cdot \mathbf{n}_i \right]^2
\]

\begin{table}[H]
\centering
\begin{tabular}{ll}
\toprule
\textbf{Paraméter} & \textbf{Végső érték} \\
\midrule
Távolság küszöb (RANSAC után) & $0{,}4 \times \text{VOXEL\_SIZE}$ \\
Maximális iteráció & 200 \\
Relatív fitness konvergencia & $10^{-6}$ \\
Relatív RMSE konvergencia & $10^{-6}$ \\
\bottomrule
\end{tabular}
\caption{ICP végső paraméterezése.}
\end{table}

% ============================================================
\section{Vizualizáció — Interaktív 3D megjelenítő}
% ============================================================

A vizualizációs rendszer teljes átalakítására került PyVista alapokra.

\subsection{Interaktív megjlenítő főbb funkciói}

Az alábbi fejlesztések történtek a vizualizáció terén:

\begin{itemize}[nosep]
    \item \textbf{Off-screen renderelés:} Batch feldolgozáshoz és automatikus riportok generálásához
    \item \textbf{Interaktív ablak:} Valós idejű forgatás, zoom, mozgatás
    \item \textbf{Fényviszonyok:} Professzionális megvilágítás és szín-kódolás
    \item \textbf{Legenda:} Automatikus szín-azonosítás az egyes pontfelhőkhöz
    \item \textbf{Vezérlési útmutató:} Beépített help az interaktív módban
\end{itemize}

% ============================================================
\section{Eredmények és Értékelés}
% ============================================================

\subsection{Saját modell — végső regisztrációs eredmények}

A fejlesztett pipeline tesztelése egy saját 3D modellon történt, amely 6 felvételből (a1\_0.ply -- a6\_0.ply) állt. Az objektum maximális dimenziója 131.6121~mm volt, amelyből a pipeline dinamikusan számított voxel méret (2.6322~mm) alapján egy 3146 pontból álló végső rekonstrukciót hozott létre.

\subsubsection{Regisztrációs eredmények}

A pipeline öt scan-t sikeresen illesztett össze az alábbi metrikákkal:

\begin{table}[H]
\centering
\begin{tabular}{lrr}
\toprule
\textbf{Scan} & \textbf{Fitness} & \textbf{RMSE (mm)} \\
\midrule
a2\_0.ply & 0.5102 & 1.330 \\
a3\_0.ply & 0.4045 & 1.370 \\
a5\_0.ply & 0.4827 & 1.302 \\
a6\_0.ply & 0.5970 & 1.279 \\
\bottomrule
\end{tabular}
\caption{Regisztrációs metrikák scan-ként.}
\end{table}

A regisztrációs folyamat során:

\begin{enumerate}[nosep]
    \item \textbf{Robusztusság:} 6 scanből 5 sikeresen illeszkedett (a4\_0.ply sérült adat miatt elvetve, majd a5 retry-val újra feldolgozva)
    \item \textbf{Regisztrációs pontosság:} Fitness értékek 0.40 -- 0.60 között, RMSE 1.3 -- 1.4 mm között
    \item \textbf{Számítási hatékonyság:} Teljes pipeline $< 5$ másodperc alatt lefutott
    \item \textbf{Dinamikus paraméterezés:} A pipeline automatikusan alkalmazkodott az objektum méretéhez (VOXEL\_SIZE = 2.6322~mm, FPFH sugár = 13.1612~mm)
\end{enumerate}

\subsubsection{Végső 3D modell karakterisztikái}

A fusion után a végső pontfelhő 3146 pontot tartalmazott, amely az eredeti 6 scan összesen 427\,980 pontjából lett előállítva ($0{,}74\%$ sűrítési arány). Az előfeldolgozási lépések során:

\begin{itemize}[nosep]
    \item Összesen 8\,550 outlier pont lett eltávolítva
    \item Voxel downsampling után 6--15\% a pontok száma
    \item Statistical Outlier Removal (SOR) 2\% -- 4\% közötti eltávolítást eredményezett
\end{itemize}

\subsection{Metrikák értelmezése}

\begin{table}[H]
\centering
\begin{tabular}{ll}
\toprule
\textbf{Metrika} & \textbf{Magyarázat} \\
\midrule
Fitness & Illeszkedő pontok aránya (0 = 0\%, 1 = 100\%) \\
Inlier RMSE & Root Mean Square Error az illeszkedő pontokon (m) \\
\#Correspondences & Az illeszkedő pontpárok száma \\
\bottomrule
\end{tabular}
\caption{Regisztrációs metrikák értelmezése.}
\end{table}

\subsection{Teljesítményprofil}

Egy tipikus regisztrációs páron mért időidőpontok:

\begin{table}[H]
\centering
\begin{tabular}{lr}
\toprule
\textbf{Szakasz} & \textbf{Idő (s)} \\
\midrule
Betöltés \& előfeldolgozás & 0.5--1.0 \\
Jellemzők kinyerése & 1.0--1.5 \\
RANSAC durva illesztés & 0.5--1.0 \\
ICP Point-to-Point & 0.2--0.5 \\
ICP Point-to-Plane & 0.2--0.5 \\
Metrikák \& exportálás & 0.2--0.3 \\
\midrule
\textbf{ÖSSZESEN} & \textbf{2.5--5.0} \\
\bottomrule
\end{tabular}
\caption{Tipikus feldolgozási idők.}
\end{table}

\section{Projekt weboldal}
% ============================================================

A projekt haladása egy dedikált weboldalon követhető, mérföldkőnként elkülönítve:

\begin{itemize}[nosep]
    \item \textbf{1.~mérföldkő:} Algoritmus terv és költségbecslés
    \item \textbf{2.~mérföldkő:} Kezdeti implementáció és első eredmények
    \item \textbf{3.~mérföldkő:} Végleges eredmények és értékelés (jelen dokumentum)
\end{itemize}

\textbf{Weboldal URL:} \url{https://ilonczai-andras.github.io/Industrial-image-processing-webpage/index.html}

% ============================================================
\begin{thebibliography}{9}
\bibitem{besl1992}
Besl, P.~J. \& McKay, N.~D. (1992).
A method for registration of 3-D shapes.
\textit{IEEE TPAMI}, 14(2), 239--256.

\bibitem{rusu2009}
Rusu, R.~B., Blodow, N. \& Beetz, M. (2009).
Fast point feature histograms (FPFH) for 3D registration.
\textit{IEEE ICRA}, 3212--3217.

\bibitem{zhou2016}
Zhou, Q.-Y., Park, J. \& Koltun, V. (2016).
Fast global registration.
\textit{ECCV}, 766--782.

\bibitem{chen1992}
Chen, Y. \& Medioni, G. (1992).
Object modelling by registration of multiple range images.
\textit{Image and Vision Computing}, 10(3), 145--155.

\bibitem{open3d2018}
Zhou, Q.-Y., Park, J. \& Koltun, V. (2018).
Open3D: A modern library for 3D data processing.
\textit{arXiv:1801.09847}.
\end{thebibliography}

\vfill
\noindent\rule{\textwidth}{0.4pt}\\
\small\textit{Ipari Képfeldolgozás Lab.\ Gyak.\ | 2.~Mérföldkő | 2026.~április}

\end{document}
